\section{Vérification exacte des constructions finies}
\label{sec:exact-reproduction}

La reproduction distingue la démonstration générale et l'évaluation d'une spécialisation. Les propositions précédentes utilisent des identités algébriques, la positivité et des hypothèses de convergence explicites. Leurs contrôles computationnels sont effectués avec des entiers et des fractions rationnelles ; le développement réticulaire conserve en outre sa représentation exacte dans l'extension algébrique déclarée. Les données du registre sont des sorties exprimées par la construction précédente. Les exécutions de cette section vérifient leur transport géométrique et opératoriel, et conservent séparément la provenance de leur génération.

Le premier contrôle restitue les séparations à partir des mineurs, vérifie la relation de Plücker dans les triangulations et compare les changements de coordinatisation aux déformations effectives. Un changement inversible des lignes multiplie tous les mineurs par le même déterminant ; les quotients marqués demeurent identiques. Un redimensionnement indépendant des colonnes conserve certains birapports et modifie, en général, les séparations reconstruites. Évaluer les deux opérations sur les mêmes données fournit un contrôle négatif de l'équivalence affirmée.

Le second contrôle subdivise les longueurs par neuf, conserve les marques héritées et vérifie les identités des moyennes et de Gram. Les déterminants sont calculés par élimination exacte et comparés à leurs produits de Vandermonde. La positivité en toute dimension finie découle de la formule démontrée ; les cas finis exécutés vérifient l'implémentation de cette formule, ses indices et ses normalisations. Les mesures de comptage et les mesures chronologiques sont maintenues séparées afin qu'une égalité de moments ne dépende pas d'un changement silencieux de normalisation.

Le troisième contrôle reproduit les huit premiers moments au moyen
de l’opérateur de Jacobi d’ordre quatre. La matrice symétrique est
représentée sous forme rationnelle dans la base des polynômes unitaires ;
sa métrique est \(\operatorname{diag}(h_0,\ldots,h_3)\). On
vérifie l’auto-adjonction dans cette métrique, la récurrence à trois
termes et l’accord entre la résolvante matricielle et sa fraction
continue en \(z=2\). Le résultat maintient explicites le vecteur
constant observé et la charge \(Q\), qui interviennent dans le
théorème de fidélité spectrale marquée.

Le quatrième groupe vérifie les changements symplectiques et leurs
hamiltoniens. La dérivée par rapport au moment récupère la
vitesse ; la transformation de Legendre est vérifiée comme
identité polynomiale ou de Laurent. L’inertie d’une forme
quadratique et la positivité d’un hamiltonien quantique
appartiennent à leurs espaces respectifs. La représentation
de Clifford est vérifiée au moyen de ses anticommutateurs,
et la composition avec l’opérateur de masse est prouvée par
le carré de l’opérateur de Dirac. Ces égalités
identifient chaque terme et ses produits croisés.

Enfin, les vérifications précédentes de triangulation, de forme canonique, de réduction de Schur, de congruence des masses et de construction réticulaire sont conservées. Dans Yang--Mills, les propositions sur la mesure commune, le raffinement et la reconstruction sont exposées dans le corps du texte avec leurs prémisses et leurs critères de réfutation. Un contrôle fini de matrices ne remplace ni le passage à la limite ni la reconstruction de champs : leurs domaines sont enregistrés séparément dans la livraison reproductible.

Le paquet de sources relie chaque programme à l’énoncé
qu’il évalue et conserve sa sortie intégrale. Les conditions sont
vérifiées par des branchements explicites, de sorte qu’elles
restent actives lors de l’exécution de l’interpréteur en mode
optimisé. Elles sont également reproduites en mode isolé, sans
charger de paquets utilisateur. La correspondance entre
construction, preuve et exécution peut être inspectée
étape par étape sans transformer la sortie d’un programme en
une assertion de portée supérieure à celle qu’il vérifie.
